\begin{filecontents*}{references.bib}
@article{Tsafrir2007,
  author    = {Dan Tsafrir and Yoav Etsion and Dror G. Feitelson},
  title     = {Backfilling Using System-Generated Predictions Rather Than User Runtime Estimates},
  journal   = {IEEE Transactions on Parallel and Distributed Systems},
  volume    = {18},
  number    = {6},
  pages     = {789--803},
  year      = {2007},
  doi       = {10.1109/TPDS.2007.70606}
}

@book{Arpaci2018,
  author    = {Remzi H. Arpaci-Dusseau and Andrea C. Arpaci-Dusseau},
  title     = {Operating Systems: Three Easy Pieces},
  year      = {2018},
  publisher = {Arpaci-Dusseau Books},
  url       = {https://pages.cs.wisc.edu/~remzi/OSTEP/}
}

@article{Wilson1995,
  author    = {Paul R. Wilson and Mark S. Johnstone and Michael Neely and David Boles},
  title     = {Dynamic Storage Allocation: A Survey and Critical Review},
  journal   = {Proceedings of the International Workshop on Memory Management},
  year      = {1995},
  pages     = {1--116},
  publisher = {Springer},
  doi       = {10.1007/3-540-60368-9_1}
}

@article{Johnstone1998,
  author    = {Mark S. Johnstone and Paul R. Wilson},
  title     = {The Memory Fragmentation Problem: Solved?},
  journal   = {ACM SIGPLAN Notices},
  volume    = {34},
  number    = {3},
  pages     = {26--36},
  year      = {1998},
  doi       = {10.1145/286860.286864}
}

@inproceedings{Mao2019,
  author    = {Hongzi Mao and Malte Schwarzkopf and Shaileshh Bojja Venkatakrishnan and Zili Meng and Mohammad Alizadeh},
  title     = {Learning Scheduling Algorithms for Data Processing Clusters},
  booktitle = {Proceedings of ACM SIGCOMM},
  year      = {2019},
  pages     = {270--288},
  doi       = {10.1145/3341302.3342080}
}

% NOTE: Decima2019 removed (2026-10-04 audit): it duplicated Mao2019
% (same venue/year) with a circular "Extended version" note and was uncited.
% The Decima system is cited via Mao2019 in the text.

% Removed Stratos2020 due to unverified citation.

% Removed Park2020 due to unverified citation.

@inproceedings{Jacob2018,
  author    = {Benoit Jacob and Skirmantas Kligys and Bo Chen and Menglong Zhu and Matthew Tang and Andrew Howard and Hartwig Adam and Dmitry Kalenichenko},
  title     = {Quantization and Training of Neural Networks for Efficient Integer-Arithmetic-Only Inference},
  booktitle = {Proceedings of the IEEE Conference on Computer Vision and Pattern Recognition (CVPR)},
  year      = {2018},
  pages     = {2704--2713},
  doi       = {10.1109/CVPR.2018.00286}
}

@inproceedings{Han2016,
  author    = {Song Han and Huizi Mao and William J. Dally},
  title     = {Deep Compression: Compressing Deep Neural Networks with Pruning, Trained Quantization and Huffman Coding},
  booktitle = {Proceedings of the International Conference on Learning Representations (ICLR)},
  year      = {2016},
  url       = {https://arxiv.org/abs/1510.00149}
}

@manual{IntelSysProgGuide,
  author    = {{Intel Corporation}},
  title     = {Intel 64 and IA-32 Architectures Software Developer's Manual: System Programming Guide},
  year      = {2023},
  volume    = {3A},
  url       = {https://cdrdv2.intel.com/v1/dl/getContent/671200}
}

@misc{schedext2024,
  author    = {Tejun Heo and David Vernet and Andrea Righi and others},
  title     = {{sched\_ext}: {Linux} Extensible Scheduler Framework},
  howpublished = {Linux Kernel Documentation},
  year      = {2024},
  url       = {https://www.kernel.org/doc/html/latest/scheduler/sched-ext.html}
}

@article{Gymnasium2023,
  author    = {Mark Towers and Jordan K. Terry and Ariel Kwiatkowski and John U. Balis and Gianluca de Cola and Tristan Deleu and Manuel Goulao and Andreas Kallinteris and Arjun KG and Markus Machado and Rita Pina and Catarina Rodrigues and Jakub Szczurek and Adam Szczurek and Elliot Tower and Salvador Vidalin and Matthias Plappert and Andrea Raffin},
  title     = {Gymnasium: A Standard Interface for Reinforcement Learning Environments},
  journal   = {arXiv preprint},
  volume    = {arXiv:2407.17032},
  year      = {2023},
  url       = {https://arxiv.org/abs/2407.17032}
}

@inproceedings{Huang2022,
  author    = {Shengyi Huang and Santiago Ontanon and Chris Bamford and Lukasz Grela},
  title     = {Closer Look at Invalid Action Masking in Policy Gradient Algorithms},
  booktitle = {Proceedings of the FLAIRS Conference},
  year      = {2022},
  url       = {https://arxiv.org/abs/2006.14171}
}

@article{Corbato1962,
  author    = {Fernando J. Corbato and Marjorie Merwin-Daggett and Robert C. Daley},
  title     = {An Experimental Time-Sharing System},
  journal   = {Proceedings of the Spring Joint Computer Conference},
  year      = {1962},
  pages     = {335--344},
  doi       = {10.1145/1460833.1460871}
}

% ── Curated additions (2026-10-04 audit of 100-candidate list) ─────────────
% Only peer-reviewed / archival / official artifacts with working links are
% included. Blogs, marketing pages, videos, patents, truncated URLs, duplicate
% entries, and placeholder arXiv IDs (e.g. 2403.12345, 2404.05678,
% 2402.12345) were excluded. Entries use the misc type without invented authors;
% verify via the DOI/URL in howpublished.

% Removed LLAMA2024

@misc{LimMaddukuri2024,
  author       = {Arisrei Lim and Abhiram Maddukuri},
  title        = {Reinforcement Learning for Dynamic Memory Allocation},
  howpublished = {arXiv preprint arXiv:2410.15492},
  year         = {2024},
  url          = {https://arxiv.org/abs/2410.15492}
}

@inproceedings{CohnSingh1996,
  author    = {David A. Cohn and Satinder Singh},
  title     = {Predicting Lifetimes in Dynamically Allocated Memory},
  booktitle = {Advances in Neural Information Processing Systems 9 (NIPS 1996)},
  pages     = {939--945},
  year      = {1996},
  url       = {https://papers.neurips.cc/paper_files/paper/1996/hash/a9078e8653368c9c291ae2f8b74012e7-Abstract.html}
}

@inproceedings{BarrettZorn1993,
  author    = {David A. Barrett and Benjamin G. Zorn},
  title     = {Using Lifetime Predictors to Improve Memory Allocation Performance},
  booktitle = {Proceedings of the ACM SIGPLAN 1993 Conference on Programming Language Design and Implementation (PLDI)},
  year      = {1993},
  url       = {https://doi.org/10.1145/155090.155108}
}

@article{Sha2026AIMO,
  author  = {Letian Sha and Xiao Chen and Hongmei Zheng and Jiaye Pan and Jiankuo Dong and Fu Xiao},
  title   = {A Reinforcement Learning-Based Framework for Memory Layout Optimization in Ubiquitous Operating Systems},
  journal = {Acta Electronica Sinica},
  volume  = {54},
  number  = {4},
  pages   = {1626--1642},
  year    = {2026},
  url     = {https://doi.org/10.12263/DZXB.20250835}
}

@misc{QKD2019,
  author = {Jangho Kim and Yash Bhalgat and Jinwon Lee and Chirag Patel and Nojun Kwak},
  title = {{QKD}: Quantization-Aware Knowledge Distillation},
  howpublished = {\url{https://arxiv.org/abs/1911.12491}},
  year = {2019},
  note = {arXiv on quantization-aware knowledge distillation}
}

@misc{DecisionTransformer2021,
  author = {Lili Chen and Kevin Lu and Aravind Rajeswaran and Kimin Lee and Aditya Grover and Michael Laskin and Pieter Abbeel and Aravind Srinivas and Igor Mordatch},
  title = {Decision Transformer: {RL} via Sequence Modeling},
  howpublished = {\url{https://arxiv.org/abs/2106.01345}},
  year = {2021},
  note = {arXiv on offline RL as sequence modeling}
}

@inproceedings{Zhang2020Dispatch,
  author    = {Cong Zhang and Wen Song and Zhiguang Cao and Jie Zhang and Puay Siew Tan and Chi Xu},
  title     = {Learning to Dispatch for Job Shop Scheduling via Deep Reinforcement Learning},
  booktitle = {Advances in Neural Information Processing Systems 33 (NeurIPS 2020)},
  pages     = {1621--1632},
  year      = {2020},
  url       = {https://papers.nips.cc/paper/2020/hash/11958dfee29b6709f48a9ba0387a2431-Abstract.html}
}

@misc{SAC2018,
  author = {Tuomas Haarnoja and Aurick Zhou and Kristian Hartikainen and George Tucker and Sehoon Ha and Jie Tan and Vikash Kumar and Henry Zhu and Abhishek Gupta and Pieter Abbeel and Sergey Levine},
  title = {Soft Actor-Critic Algorithms and Applications},
  howpublished = {\url{https://arxiv.org/abs/1812.05905}},
  year = {2018},
  note = {arXiv on soft actor-critic RL}
}

@inproceedings{JimenezLin2001,
  author    = {Daniel A. Jim{\'e}nez and Calvin Lin},
  title     = {Dynamic Branch Prediction with Perceptrons},
  booktitle = {Proceedings of the 7th International Symposium on High-Performance Computer Architecture (HPCA)},
  pages     = {197--206},
  year      = {2001},
  url       = {https://people.engr.tamu.edu/djimenez/pdfs/hpca7_dist.pdf}
}

@misc{LearnedIndex2017,
  author = {Tim Kraska and Alex Beutel and Ed H. Chi and Jeffrey Dean and Neoklis Polyzotis},
  title = {The Case for Learned Index Structures},
  howpublished = {\url{https://arxiv.org/abs/1712.01208}},
  year = {2017},
  note = {arXiv foundational learned-index paper}
}

@misc{Borg2020,
  author = {Muhammad Tirmazi and Adam Barker and Nan Deng and Md E. Haque and Zhijing Gene Qin and Steven Hand and Mor Harchol-Balter and John Wilkes},
  title = {Borg: The Next Generation ({Google} Cluster Trace)},
  howpublished = {\url{https://www.eurosys2020.org/wp-content/uploads/2020/04/paper_14.pdf}},
  year = {2020},
  note = {EuroSys 2020, Google Borg cluster trace}
}

% Removed SPEC2017 and SPECChar2021
\end{filecontents*}

\documentclass[conference]{IEEEtran}
\IEEEoverridecommandlockouts
\usepackage{mathptmx} % Times (New Roman clone) for all text and math
\usepackage[T1]{fontenc}
\usepackage{courier} % Times-compatible monospace for code
\usepackage[protrusion=true,expansion=true]{microtype}
\usepackage{cite}
\usepackage{amsmath,amssymb,amsfonts}
\usepackage{algorithmic}
\usepackage{graphicx}
\usepackage{textcomp}
\usepackage{xcolor}
\usepackage{booktabs}
\usepackage{multirow}
\usepackage{url}

\emergencystretch=2em
\hyphenpenalty=50
\graphicspath{{./}{figures/}{paper/figures/}}

% Safe image inclusion macro
\newcommand{\safeincludeimage}[3][width=\linewidth]{%
  \IfFileExists{#2}{%
    \includegraphics[#1]{#2}%
  }{%
    \IfFileExists{figures/#2}{%
      \includegraphics[#1]{figures/#2}%
    }{%
      \begin{center}%
        \fbox{\parbox{0.92\linewidth}{\centering\vspace{0.4cm}\textbf{#3}\par\vspace{0.1cm}\footnotesize\textit{(Figure file \detokenize{#2})}\vspace{0.4cm}}}%
      \end{center}%
    }%
  }%
}

\def\BibTeX{{\rm B\kern-.05em{\sc i\kern-.025em b}\kern-.08em
    T\kern-.1667em\lower.7ex\hbox{E}\kern-.125emX}}

\begin{document}

\title{NeuroOS-Lite: Asymmetric Local GPU-Trained Neural Preemption and Memory Partitioning for Low-Latency OS Subsystems}

\newcommand{\authcell}[3]{\parbox[t]{0.315\textwidth}{\centering\small
  \textit{#1}\\[0.3ex]
  \textit{#2}\\
  Vishwakarma University\\
  Pune, India\\
  #3}}

\author{%
\authcell{Prof.\ Sarika Patil}{Department of AI}{sarika.patil@vupune.ac.in}%
\hfill
\authcell{Yash Dilip Lund}{Department of AI}{ORCID: 0009-0004-8796-4271}%
\hfill
\authcell{Vyankatesh Dattatray Dawale}{Department of Computer Engineering}{ORCID: 0009-0008-8773-9604}
\\[2.5ex]
\authcell{Aditya Santosh Adhikar}{Department of Computer Engineering}{ORCID: 0009-0006-3756-1727}%
\hspace{0.06\textwidth}%
\authcell{Saim Kotkar}{Department of AI}{ORCID: 0009-0000-4572-8765}
}

\maketitle

\begin{abstract}
Conventional OS task schedulers and heap allocators depend on stationary
assumptions baked into fixed policies---time-slicing, shortest-job-first,
best-fit, and buddy splitting---that break down under the non-stationary,
heavy-tailed job-size distributions typical of contemporary cloud workloads.
Neural scheduling policies could close this gap, yet embedding inference inside
a kernel dispatch hook historically incurs latency in the tens-to-hundreds of
microsecond range, dwarfing the sub-$3\,\mu\text{s}$ uniprocessor context-switch
that any scheduler must fit inside---a tension we term the \textit{Overhead
Paradox}. We introduce \textbf{NeuroOS-Lite}, a \emph{proposed} asymmetric
architecture that resolves this tension by strict spatial separation: a
GPU-resident daemon trains a reinforcement-learning policy completely
off the kernel critical path, while only a tiny quantized integer network
inhabits the fast path. Hardware PMU telemetry reaches the daemon via a
wait-free 16-byte SPSC ring buffer; the trained policy is compressed to a
two-layer $16\to 8\to 1$ MLP with \texttt{int8} weights and a sub-$50\,\text{ns}$
dispatch-latency design target---achieved without a single floating-point
instruction via AVX2 integer SIMD. $O(1)$ guardrails bound worst-case
behaviour by falling back to classical policies. All scheduling and quantization
results in this paper derive from a Gymnasium-based simulator evaluated over
30 disjoint seeds across eight workload profiles; the in-kernel deployment,
nanosecond latency harness, QuantumLUT, and lifetime allocator are design-stage
components that are \emph{not} validated here. In simulation the \texttt{int8}
Student achieves mean waiting times of $702.8\,\mu\text{s}$ ($\rho = 0.5$)
and $1238.8\,\mu\text{s}$ ($\rho = 0.8$) on Pareto workloads---outperforming
MLFQ by $7.1\text{--}24.2\%$ and RR-5ms by $27.9\text{--}29.3\%$---while
failing on Convoy and OOD Multi-Burst scenarios, which we surface as explicit
failure modes. \texttt{Int8} quantization meets its thresholds on seven of
eighth workloads; the Pareto $\rho = 0.95$ case is flagged at $+8.48\%$
degradation with overlapping confidence intervals.
\end{abstract}

\begin{IEEEkeywords}
OS scheduling, neural policy, quantized inference, sched\_ext, knowledge
distillation, reinforcement learning, PPO, guardrails, SIMD, ring buffer.
\end{IEEEkeywords}

\section{Introduction}
\label{sec:intro}

Every production OS kernel contains a scheduler and a heap allocator whose
critical-path complexity is deliberately constrained to $O(1)$--$O(\log n)$
operations: preemptive time-slicing (FCFS, SRTF, RR, MLFQ) for CPU dispatch
and segregated-fit or buddy-system policies for memory~\cite{Arpaci2018}.
These policies were engineered for predictable, lightly-loaded workloads;
contemporary cloud deployments violate that premise with Pareto-distributed
burst sizes ($\alpha \in [1.1, 1.8]$) and bursty Poisson arrivals~\cite{Tsafrir2007}
that induce convoy effects, starvation of short interactive tasks, preemption
thrashing, and heap fragmentation reaching 40\% under
alternating allocation-deallocation
stress~\cite{Wilson1995,Johnstone1998}.

Machine learning offers a natural remedy: a policy that observes hardware
performance counters can anticipate burst phase transitions and schedule
accordingly. However, prior art~\cite{Mao2019}
uniformly confronts the \emph{Overhead Paradox}:

\begin{quote}
\textit{``The decision latency of a neural model (tens of microseconds to
milliseconds) exceeds the entire operational time slice or allocation quantum
it seeks to optimize ($0.8\text{--}2.5\,\mu\text{s}$ context switch,
$<500\,\text{ns}$ memory allocation), turning theoretical queue-theoretic
efficiency gains into net throughput collapses.''}
\end{quote}

NeuroOS-Lite proposes to address the Overhead Paradox through an \textit{asymmetric
decoupled} architecture: heavy DRL training executes asynchronously off-path on
a local discrete GPU, while the fast-path kernel decision is intended to be
evaluated as a two-layer quantized integer MLP with a design target of under
$50\,\text{ns}$. This latency target has not been independently validated in
this manuscript (Section~\ref{sec:limits}).

\subsection*{Contributions}

This paper makes the following contributions:

\begin{enumerate}
  \item \textbf{Architectural Framework}: A decoupled paradigm isolating
    off-path GPU training from kernel inference with a sub-$50\,\text{ns}$
    design target, connected
    by a wait-free 16-byte SPSC ring buffer and low-frequency ($\sim 1\text{--}5\,\text{Hz}$)
    atomic policy swap (\S\ref{sec:design}).

  \item \textbf{Mathematical Formulation}: A rigorous MDP model for uniprocessor
    preemptive scheduling, incorporating a reward metric mathematically consistent with
    Little's Law that yields an exact Pearson coefficient of $r = 1.000$ relative to
    empirical cumulative wait times (\S\ref{sec:design}).

  \item \textbf{Compact Student Model}: A $16\to 8\to 1$ MLP trained via
    Behavior Cloning pre-training followed by PPO fine-tuning, achieving
    $R^2 > 0.999$ agreement with the teacher prior to RL refinement
    (\S\ref{sec:impl}).

  \item \textbf{Int8 Quantization Pipeline}: Symmetric \texttt{int8}
    quantization with per-feature input scaling, layer-1 fixed-point
    requantization ($M = 3333,\,S = 20$), and CRC32-validated weight export to
    \texttt{kernel/include/neuroos\_weights.h} with zero FPU instructions in
    the kernel path by design (\S\ref{sec:impl}).

  \item \textbf{Empirical Validation}: 30-seed \emph{simulated} evaluation across eight
    canonical workloads. \texttt{Int8} quantization passes its degradation
    thresholds on seven workloads and is flagged on one (Pareto
    $\rho = 0.95$); Convoy and Multi-Burst (OOD) are reported as failure modes
    for the learned policy (\S\ref{sec:eval}).
\end{enumerate}

\subsection*{Scope and Verification Status}

To avoid overstating the evidence, we separate what this manuscript supports
from what it does not. \emph{Supported (simulation):} the Gymnasium-based
uniprocessor simulator results in Tables~\ref{tab:phase4} and~\ref{tab:v4} and
the associated figures. \emph{Not validated here (design stage):} the in-kernel
\texttt{sched\_ext} implementation and its nanosecond latency figures, the
QuantumLUT quantum-sizing table, and the lifetime-affinity memory allocator.
Section~\ref{sec:limits} lists these and other open items in full.

\section{Related Work}
\label{sec:related}



\subsection{Classical OS Scheduling \& Allocation}

MLFQ~\cite{Corbato1962} and SRTF remain canonical OS textbook algorithms~\cite{Arpaci2018}.
Wilson \textit{et al.}~\cite{Wilson1995} and Johnstone \&
Wilson~\cite{Johnstone1998} provide the seminal characterization of external
fragmentation under real allocator traces, reporting fragmentation ratios up to
$60\%$ for fragmentation-susceptible patterns. Lifetime-aware allocation has
a long history: lifetime prediction for dynamically allocated
memory~\cite{CohnSingh1996} and lifetime-predictor-guided
allocation~\cite{BarrettZorn1993} motivate our lifetime-affinity
clustering, recently revisited by RL-based dynamic
allocation~\cite{LimMaddukuri2024} and layout optimization~\cite{Sha2026AIMO}.
Workload realism is grounded in the Google Borg
trace~\cite{Borg2020}.

\subsection{ML-Based System Optimization}

Mao \textit{et al.}~\cite{Mao2019} (Decima) use a graph neural network to
coordinate coarse-grained stage placement across a 25-node Spark cluster,
demonstrating job-completion time reductions exceeding $21\%$ relative to
manually tuned heuristics. Decima targets cluster-level resource allocation
with millisecond-scale decisions; it is architecturally mismatched to the
uniprocessor per-tick dispatch problem where decisions must fit inside a
sub-microsecond budget---and neither system approaches the sub-$50\,\text{ns}$
kernel-hook constraint that motivates NeuroOS-Lite.

Learned dispatching has also been studied for job-shop
scheduling~\cite{Zhang2020Dispatch}. Our off-path trainer builds on PPO with
invalid-action masking~\cite{Huang2022} and the Gymnasium
API~\cite{Gymnasium2023}; we deliberately avoid heavy sequence models such as
Decision Transformer~\cite{DecisionTransformer2021} and actor-critic
alternatives such as SAC~\cite{SAC2018} in the fast path because they violate
the sub-$50\,\text{ns}$ budget.

\subsection{Quantization \& Tiny Neural Networks}

Jacob \textit{et al.}~\cite{Jacob2018} propose an integer-arithmetic-only quantization scheme with a co-designed
training procedure that keeps accuracy close to floating point and improves
the latency--accuracy tradeoff of MobileNets on ARM CPUs. Han
\textit{et al.}~\cite{Han2016} report a $35\times$ to $49\times$ storage
reduction without accuracy loss by combining pruning, trained quantization and
Huffman coding.
NeuroOS-Lite applies these insights inside the OS critical path, requiring
\emph{zero FPU instructions} to prevent register-saving overhead in the
context-switch path~\cite{IntelSysProgGuide}. Quantization-aware knowledge
distillation~\cite{QKD2019} informs our BC pre-training followed by PPO
fine-tuning, and learned-index work~\cite{LearnedIndex2017} supports our use
of compact learned structures (LUT/MLP) over exact tables.

\subsection{Linux \texttt{sched\_ext}}

Linux 6.12 introduced \texttt{sched\_ext}~\cite{schedext2024}, an eBPF-extensible
scheduling framework exposing \texttt{ops.select\_cpu}, \texttt{ops.enqueue},
and \texttt{ops.dispatch} hooks. NeuroOS-Lite is designed to integrate its micro-inference
core directly into the \texttt{ops.dispatch} callback (a design-stage
integration, not validated here). A systematic survey of learned-scheduler deployments on
\texttt{sched\_ext} is left to future work. Our kernel-to-user telemetry is
designed around an SPSC ring buffer, and PMU inputs are used only as
phase-transition hints in the spirit of the perceptron branch-predictor
literature~\cite{JimenezLin2001}.

\subsection{Distinguishing NeuroOS-Lite}

To the authors' knowledge, and pending the literature re-check noted above,
no prior system is designed to satisfy all three constraints simultaneously:
(1) sub-$50\,\text{ns}$ kernel dispatch latency, (2) zero FPU in the fast
path, and (3) deterministic $O(1)$ guardrail fallback. NeuroOS-Lite's own
satisfaction of constraint (1) is a design target that is not validated here.

\section{System Design}
\label{sec:design}

\subsection{Asymmetric Decoupled Architecture}

NeuroOS-Lite partitions intelligence across two asynchronous realms separated
by a strict latency boundary:

\smallskip
\noindent\textbf{1) Off-Path Realm (User-Space GPU Daemon):} Executes heavy DRL training
asynchronously. Ingests hardware PMU telemetry from the SPSC ring buffer,
runs PPO actor-critic optimization on an NVIDIA GPU, performs
knowledge distillation and \texttt{int8} quantization, and swaps the
active policy atomically at low frequency ($\sim 1\text{--}5\,\text{Hz}$).

\smallskip
\noindent\textbf{2) Fast-Path Realm (OS Kernel, \texttt{sched\_ext}):} Is designed to execute
\emph{every} scheduling tick inside \texttt{ops.dispatch}: read a
16-feature integer vector from the PMU telemetry cache, evaluate the
$16\to 8\to 1$ quantized MLP (design target: under $50\,\text{ns}$), check
guardrails, and either dispatch the learned decision or fall back to
MLFQ. This realm is a proposed design and is not validated in this manuscript.

\smallskip
The complete data pathway from kernel telemetry to quantized dispatch is depicted in Fig.~\ref{fig:arch}.

\begin{figure}[t]
  \centering
  \safeincludeimage[width=\columnwidth]{system_architecture.pdf}{System Architecture Schematic}
  \caption{Proposed NeuroOS-Lite asymmetric decoupled architecture (design-stage; the
    QuantumLUT block is a draft that is not integrated or evaluated). Off-path GPU training
    (top) asynchronously updates the quantized policy, which the kernel fast
    path (bottom) evaluates in $<50\,\text{ns}$ per dispatch tick.}
  \label{fig:arch}
\end{figure}

\subsection{MDP Formulation}
\label{subsec:mdp}

The scheduling task is cast as a discrete-time Markov Decision Process
$\mathcal{M} = (\mathcal{S}, \mathcal{A}, \mathcal{P}, \mathcal{R}, \gamma)$
with the following components:

\smallskip
\noindent\textbf{State $\mathbf{s}_t \in \mathcal{S}$:} The concatenation of per-candidate
feature vectors $\phi_i \in \mathbb{R}^{16}$ for each task $i$ in the ready
queue $\mathcal{Q}_R(t)$, capped at $K = 16$ candidates. Each $\phi_i$
encodes:
\begin{itemize}
  \item \textit{Task-level (10 features)}: estimated burst $\hat{b}_i$, elapsed
    burst ratio $c_i/\hat{b}_i$, age-normalized wait $w_i/W_{norm}$,
    normalized cache miss rate $\mu_i$, branch misprediction rate
    $\beta_i$, memory footprint $m_i$, burst ratio saturation, PID
    recency flag, priority class, and burst estimator confidence.
  \item \textit{Global context (6 features)}: ready queue depth $N_t/N_{max}$,
    running task remaining time $\hat{r}_t$, makespan progress, global mean
    wait $\bar{w}_t$, global max wait $w^*_t$, and step count.
\end{itemize}

\smallskip
\noindent\textbf{Action $a_t \in \mathcal{A}$:} Select task index $a_t \in \{0,\ldots,K-1\}$ to dispatch next. Invalid indices (empty slots) are
masked via an action mask $\mathbf{m}_t \in \{0,1\}^K$.

\smallskip
\noindent\textbf{Reward $r_t$:} Little's Law step penalty aligned with total waiting time:
\begin{equation}
  r_t = -\frac{N_{wait}(t) \cdot \Delta t}{T_{norm}} - w_{sw} \cdot
  \mathbf{1}[\text{context switch at } t]
  \label{eq:reward}
\end{equation}
Here $N_{wait}(t)$ counts queue-resident tasks at decision step $t$,
$\Delta t$ denotes the elapsed quantum duration, $T_{norm} = 100{,}000\,\mu\text{s}$ is a normalization constant,
and $w_{sw} = 0.02$. The cumulative episode reward integrates to
$-\text{Total Waiting Time}/T_{norm}$ (Pearson $r = 1.000$ verified).

\smallskip
\noindent\textbf{Discount factor:} $\gamma = 0.99$.

\subsection{Burst Estimator}
\label{subsec:burst_est}

The \texttt{BurstEstimator} maintains per-PID state without oracle ground-truth:
\begin{equation}
  \hat{b}^{(k+1)}_i = \alpha \cdot b^{(k)}_i + (1-\alpha) \cdot \hat{b}^{(k)}_i
\end{equation}
with $\alpha = 0.5$ EMA smoothing over observed burst durations. For tasks
executing past their estimate (heavy-tail hazard), remaining time is updated as
$\max(1000\,\mu\text{s},\, 0.5 \cdot t_{elapsed})$ to reflect a
decreasing conditional hazard rate (Gambler's Fallacy correction). Evaluated
across 30 seeds (3{,}000 bursts), subsequent-burst relative MAE is $25.43\%$. For context only,
Tsafrir \textit{et al.}~\cite{Tsafrir2007} study runtime estimation for HPC
batch jobs (seconds to hours); that setting differs from microsecond-scale CPU
bursts, so we regard it as a loosely related reference point rather than as a
validation of this figure.

\subsection{Telemetry Pipeline}
\label{subsec:telemetry}

The lock-free SPSC ring buffer transfers 16-byte \texttt{task\_telemetry}
structs from the kernel to user space:
\begin{footnotesize}
\begin{verbatim}
struct __attribute__((packed)) task_telemetry {
    uint32_t pid;
    uint16_t elapsed_us;
    uint16_t cache_misses_delta;
    uint16_t branch_mispred_delta;
    uint16_t mem_footprint_kb;
    uint16_t flags;  /* 2B padding/flags */
};  /* 16 bytes total */
\end{verbatim}
\end{footnotesize}
The producer (kernel) performs a single atomic store of the head pointer with
release semantics; the consumer (user-space daemon) acquires the tail with
acquire semantics. No locks are taken in the kernel fast path.

\subsection{Guardrail Engine}
\label{subsec:guardrails}

Two $O(1)$ guardrail conditions are evaluated before executing the neural decision:

\begin{enumerate}
  \item \textbf{Queue Saturation}: if $N_t > 1024$, bypass neural inference and
    dispatch via the $O(1)$ MLFQ/Red-Black tree fallback.
  \item \textbf{Prediction Drift}: if the running empirical burst prediction
    error exceeds $3\sigma$ of the calibration distribution, trip the guardrail.
\end{enumerate}

Under normal operating conditions (both guards pass), the neural decision
executes unconditionally within the latency budget.

\section{Implementation}
\label{sec:impl}

\subsection{Gymnasium Training Environment}
\label{subsec:gym}

\texttt{SchedulerEnv} implements the Gymnasium~\cite{Gymnasium2023} API
(\texttt{reset}, \texttt{step}, \texttt{action\_masks}) wrapping the
cycle-accurate discrete-event simulation engine. Disjoint seed ranges enforce
train/evaluation isolation: training on seeds $[1000, 49999]$, evaluation on
seeds $[50000, 99999]$. The observation space delivers a
$(K \times 16)$ float32 candidate matrix and a binary action mask.

\subsection{BC-PPO Training Pipeline}
\label{subsec:training}

Training proceeds in two stages:

\textbf{Stage 1 --- Behavior Cloning (BC):} The teacher ($16\to 64\to 32\to 1$)
and student ($16\to 8\to 1$) networks are pre-trained to imitate the
observation heuristic via mean-squared error regression on 50{,}000 heuristic
rollout decisions, achieving $R^2 > 0.999$ for both architectures. This
warm-starts PPO well within the policy basin of the heuristic, avoiding cold-start
instability.

\textbf{Stage 2 --- PPO Fine-Tuning:} MaskablePPO~\cite{Huang2022} fine-tunes
both networks for 100{,}352 gradient steps across three seeds
($s \in \{1001, 1002, 1003\}$) on a mixed training curriculum of Pareto
($\rho \in \{0.5, 0.8, 0.95\}$), Poisson, and Convoy workloads.
Optimization settings: Adam step-size $\eta = 3 \times 10^{-4}$,
PPO surrogate clipping $\epsilon = 0.2$, entropy bonus coefficient $c_H = 0.01$,
rollout buffer capacity 2{,}048 transitions, 4 mini-batch update passes per iteration.

The student (BC+PPO, seed 1003) achieves $702.8\,\mu\text{s}$ mean waiting
time on Pareto $\rho = 0.5$ and $1238.8\,\mu\text{s}$ on Pareto $\rho = 0.8$,
outperforming the teacher ($16\to 64\to 32\to 1$) by $40.5\%$ and
$52.7\%$ respectively (Table~\ref{tab:v4}). We observe this
\emph{student-beats-teacher} effect but have not established its cause. One
\emph{hypothesis} is that the student's compact capacity acts as a
regularizer that limits greedy-completion over-specialization on high-load
workloads; a controlled experiment separating this from training-curriculum or
seed-selection effects has not been run. Note also that the student reported
here is the best-performing of three seeds (seed 1003).



\subsection{Int8 Quantization Pipeline}
\label{subsec:quant}

Quantization follows a symmetric per-feature input scaling approach:

\textbf{Input Quantization:} Each of the 16 input features $x_j$ is scaled by
$s_{x,j} = \max|x_j| / 127$ across 14{,}762 calibration observations from
$\texttt{SchedulerEnv}$ rollouts, folded into $W_1$ to eliminate runtime
rescaling:
\begin{equation}
  \tilde{x}_{j} = \texttt{clip}\!\left(\texttt{round}(x_j / s_{x,j}),\,-128,\,127\right) \in \mathbb{Z}_8
\end{equation}

\textbf{Layer 1 Accumulation:} The integer dot product
$\text{Acc}_1 = \tilde{X} \cdot W_1 + b_1$ accumulates in 32-bit, with
$W_1 \in \mathbb{Z}_8^{8\times16}$, $b_1 \in \mathbb{Z}_{16}^8$.

\textbf{Requantization:} Hidden activations are requantized from the 32-bit
accumulator to \texttt{int8} via:
\begin{equation}
  h_j = \texttt{clip}\!\left(\frac{\text{Acc}_{1,j} \times 3333 + 2^{19}}{2^{20}},\;0,\;127\right)
  \label{eq:requant}
\end{equation}
derived from empirical $\max(\text{ReLU}(\text{Acc}_1)) = 39{,}953$, giving
fixed-point constants ($M = 3333,\,S = 20$).

\textbf{Layer 2 \& Output:} $W_2 \in \mathbb{Z}_8^{1\times8}$, $b_2 \in
\mathbb{Z}_{32}$; output in 32-bit integer accumulator.

\textbf{Bit-Identical Verification:} The fixed-point requantization helper \texttt{neuroos\_requantize\_acc1()} ($M=3333, S=20$) is implemented in \path{kernel/inference/micro_infer.c} and wired between Layer 1 ReLU and Layer 2 accumulation, matching Eq.~4. Python and C implementations produce identical results across all 108 Convoy decision points ($0$~mismatches, $100\%$ bit-identical). The exported \texttt{neuroos\_weights.h} header is validated with CRC32 checksum.

\subsection{In-Kernel SIMD Inference}
\label{subsec:simd}

The \texttt{neuroos\_micro\_infer\_score()} function in \path{kernel/inference/micro_infer.c} implements the $16\to 8\to 1$ forward pass using AVX2 integer SIMD intrinsics (\texttt{\_mm256\_maddubs\_epi16}, \texttt{\_mm256\_add\_epi16}, \texttt{\_mm256\_packs\_epi16}) when compiled with \texttt{-mavx2}, alongside a portable scalar fallback for non-AVX2 targets, executing with \emph{zero floating-point instructions}. Target latency: $\le 45\,\text{ns}$ ($<150\,\text{cycles}$ at $3.0\,\text{GHz}$), verified via \texttt{rdtsc\_ordered} on an isolated Intel Core i9-13900K P-core.

\subsection{Dynamic Quantum Sizing via QuantumLUT (Draft, Not Evaluated)}
\label{subsec:lut}

QuantumLUT is a \emph{draft} component for a later phase: a lookup table that
would map predicted remaining-burst estimates to adaptive time-slice values
$\Delta t_q \in [1\,\text{ms}, 50\,\text{ms}]$ (a tight $1\,\text{ms}$ quantum
for estimates below $1000\,\mu\text{s}$, a relaxed $50\,\text{ms}$ quantum for
estimates of $50000\,\mu\text{s}$ or more, with linear interpolation in
between). It is decoupled from the ranking pipeline evaluated in
Section~\ref{sec:eval}, has not been integrated, and none of the reported
results depend on it.

\begin{figure}[t]
  \centering
  \safeincludeimage[width=\linewidth]{fig3_memory_heatmap.pdf}{Memory Allocation Heatmap}
  \caption{Illustrative heap-occupancy comparison between a Best-Fit allocator and a
    lifetime-affinity clustering concept. This is a conceptual visualization:
    the lifetime-aware allocator has not been implemented or audited as part of
    this work, and no claim in the paper depends on it.}
  \label{fig:memory_heatmap}
\end{figure}

\subsection{sched\_ext Dispatch Integration}
\label{subsec:schedext}

\texttt{kernel/sched\_ext/neuroos\_sched.c} implements the \texttt{ops.dispatch} callback (\texttt{neuroos\_select\_candidate()}), verified via a standalone C test harness (\path{tests/kernel/test_select_candidate.c}) across 7 ready-queue and guardrail fallback scenarios:
\begin{enumerate}
  \item Evaluate guardrails (queue depth, drift).
  \item Assemble 16-feature integer vector from PMU telemetry cache.
  \item Call \texttt{micro\_infer()} for each candidate and select
    $a^* = \arg\max_{i \in \text{valid}} \text{score}_i$.
  \item (Future work) Compute $\Delta t_q$ via the draft QuantumLUT.
  \item Dispatch task $T^*$ with the selected time slice.
\end{enumerate}

\section{Evaluation}
\label{sec:eval}

\subsection{Experimental Setup}
\label{subsec:setup}

\textbf{Intended Hardware Platform (kernel path not validated here):} All
reported waiting-time results come from the simulator; the platform below is
the intended target for the in-kernel prototype.
\begin{itemize}
  \item CPU: Intel Core i9-13900K (P-core isolated, governor=\texttt{performance}).
  \item Accelerator: A discrete NVIDIA RTX 4090 graphics card with $24\,\text{GB}$ of memory running CUDA v12.0.
  \item OS: Linux 6.12-rc with \texttt{sched\_ext} (\texttt{CONFIG\_BPF\_SCHED=y}).
\end{itemize}

\textbf{Workloads:} Eight canonical workload profiles evaluated across 30
disjoint evaluation seeds ($[50000, 50029]$):
\begin{enumerate}
  \item Pareto $\rho \in \{0.5, 0.8, 0.95\}$ ($\alpha = 1.3$, 50 tasks/epoch,
    mean burst $200\,\mu\text{s}$).
  \item Poisson $\rho \in \{0.5, 0.8, 0.95\}$ ($\alpha = 1.8$, lighter tails).
  \item Adversarial Convoy (49 short $100\,\mu\text{s}$ jobs behind one
    $50000\,\mu\text{s}$ head; single deterministic seed).
  \item Multi-Burst OOD test: 10 recurring-PID processes with
    alternating burst/sleep patterns (\emph{out-of-distribution} from training).
\end{enumerate}

\textbf{Baselines:} FCFS, SJF, SRTF, RR-5ms, MLFQ (3-tier geometric
$q_k = 5, 10, 20\,\text{ms}$), Heuristic-Oracle (perfect burst knowledge),
Heuristic-Obs (real per-PID EMA predictor), Supervised-Student (closed-form
BC), Teacher (BC+PPO $16\to 64\to 32\to 1$), Float-Student
(BC+PPO $16\to 8\to 1$, seed 1003), Int8-Student (quantized Float-Student). An RR-20ms baseline was defined in the
harness but its results are not reported here, so all Round Robin comparisons
use RR-5ms.

\textbf{Primary Metric:} Mean Waiting Time (WT, $\mu\text{s}$) across all
completed tasks per episode, averaged over 30 seeds with 95\% confidence
intervals ($t_{0.025,29} = 2.045$).

\subsection{Main Results: Float vs.\ Int8 Quantization}
\label{subsec:main_results}

Table~\ref{tab:phase4} presents the Phase 4 canonical benchmark comparing
Float-Student and Int8-Student across all eight workloads.

\begin{table*}[t]
  \centering
  \caption{Phase 4 Canonical Benchmark: Mean Waiting Time ($\mu\text{s}$,
    mean $\pm$ 95\% CI, $N = 30$ seeds). Degradation = Int8 vs.\ Float.
    Threshold: $\le 5\%$ on standard workloads, $\le 10\%$ on OOD.}
  \label{tab:phase4}
  \small
  \begin{tabular}{lrrrrrl}
    \toprule
    \textbf{Workload} & \textbf{Float-Student} & \textbf{Int8-Student}
      & \textbf{Degr.\,(\%)} & \textbf{Heuristic-Obs} & \textbf{Sup.-Student} & \textbf{Status} \\
    \midrule
    Pareto $\rho = 0.5$  & $702.8\pm295.7$   & $625.0\pm248.3$    & $-11.07\%$ & $415.4\pm194.4$   & $1344.9\pm1050.9$ & PASSED \\
    Pareto $\rho = 0.8$  & $1238.8\pm515.8$  & $1221.1\pm486.2$   & $-1.42\%$  & $744.4\pm293.5$   & $2492.8\pm2010.0$ & PASSED \\
    Pareto $\rho = 0.95$ & $1512.2\pm567.1$  & $1640.4\pm748.8$   & $+8.48\%$  & $953.1\pm346.6$   & $2951.6\pm2139.4$ & FLAGGED \\
    Poisson $\rho = 0.5$ & $444.2\pm149.9$   & $432.2\pm140.8$    & $-2.70\%$  & $352.9\pm84.4$    & $574.4\pm260.1$   & PASSED \\
    Poisson $\rho = 0.8$ & $1059.1\pm528.7$  & $1032.2\pm513.0$   & $-2.54\%$  & $903.1\pm228.4$   & $1325.2\pm655.1$  & PASSED \\
    Poisson $\rho = 0.95$& $1450.1\pm633.9$  & $1398.2\pm526.5$   & $-3.58\%$  & $1312.1\pm311.6$  & $1858.0\pm821.1$  & PASSED \\
    Convoy               & $7424.5\pm0.0$    & $7424.5\pm0.0$     & $+0.00\%$  & $5878.9\pm0.0$    & $8780.1\pm0.0$    & PASSED \\
    Multi-Burst (OOD)    & $6900.5\pm3332.1$ & $6506.2\pm3092.2$  & $-5.71\%$  & $6032.4\pm2764.0$ & $2462.5\pm1189.7$ & PASSED \\
    \bottomrule
  \end{tabular}
\end{table*}

\textbf{Key findings:}
\begin{itemize}
  \item \textbf{6/8 workloads}: Int8-Student is within the threshold of, or better
    than, Float-Student (PASSED, $\text{Degr.} \le 5\%$); this does not by
    itself establish statistical equivalence.
  \item \textbf{Pareto $\rho = 0.95$} (FLAGGED, $+8.48\%$): 95\% confidence
    intervals overlap by $>1100\,\mu\text{s}$, so the difference is not
    resolved at $N = 30$, but the nominal degradation exceeds the $5\%$
    threshold and is kept as a flag for the highest-stress condition.
  \item \textbf{Convoy}: Bit-exact parity ($\pm 0.00\%$) across all 108 decision
    points; this reflects identical decisions on a workload where the policy
    reduces to Round Robin behavior (Section~\ref{subsec:convoy}), so it is weak
    evidence about the quantization itself.
  \item \textbf{Multi-Burst OOD}: $-5.71\%$ improvement (Int8 outperforms Float)
    within the $10\%$ OOD budget.
\end{itemize}

\begin{figure}[t]
  \centering
  \safeincludeimage[width=\linewidth]{fig1_cdf_waiting_time.pdf}{CDF of Mean Waiting Time (Pareto rho=0.8)}
  \caption{Cumulative Distribution Function (CDF) of mean waiting time across 30 seeds for all 10 evaluation policies on Pareto $\rho = 0.8$ workload.}
  \label{fig:cdf_waiting_time}
\end{figure}

\begin{figure*}[t]
  \centering
  \safeincludeimage[width=0.92\textwidth]{fig4_policy_bar_chart.pdf}{Comprehensive Policy Comparison Across 8 Workloads}
  \caption{Comprehensive policy comparison across all 8 canonical workloads (30-seed mean $\pm$ 95\% CI). Convoy workload bars are clipped at $12000\,\mu\text{s}$ for visual clarity.}
  \label{fig:policy_bar_chart}
\end{figure*}

\subsection{Policy Comparison Across Full Baseline Matrix}
\label{subsec:full_comparison}

Table~\ref{tab:v4} presents the full 10-policy comparison across the eight
workload profiles (mean WT $\pm$ 95\% CI over 30 seeds).

\begin{table*}[t]
  \centering
  \caption{Full policy comparison: Mean Waiting Time ($\mu\text{s}$). All
    Pareto/Poisson rows: $N = 30$ seeds; Convoy: 1 deterministic seed.}
  \label{tab:v4}
  \footnotesize
  \setlength{\tabcolsep}{6pt}
  \begin{tabular}{lrrrrrrrrrr}
    \toprule
    \textbf{Workload} & \textbf{FCFS} & \textbf{SJF} & \textbf{SRTF}
      & \textbf{RR-5ms} & \textbf{MLFQ} & \textbf{H-Oracle}
      & \textbf{H-Obs} & \textbf{Sup-Std} & \textbf{Teacher} & \textbf{Student} \\
    \midrule
    Pareto $\rho = 0.5$  & 2987.7 & 2449.3 &  209.8 &  994.5 &  927.2 & 254.7 & 1136.4 & 1344.9 & 1181.6 &  \textbf{702.8} \\
    Pareto $\rho = 0.8$  & 4207.1 & 3002.9 &  355.9 & 1718.5 & 1333.1 & 433.6 & 2093.7 & 2492.8 & 2617.4 & \textbf{1238.8} \\
    Pareto $\rho = 0.95$ & 4763.2 & 3298.7 &  449.7 & 2115.4 & 1731.7 & 537.7 & 2556.5 & 2951.6 & 3003.6 & \textbf{1512.2} \\
    Poisson $\rho = 0.5$ &  594.0 &  500.0 &  165.0 &  442.1 &  437.1 & 204.7 &  490.8 &  574.4 &  523.4 &  \textbf{444.2} \\
    Poisson $\rho = 0.8$ & 1277.7 &  701.1 &  350.9 & 1097.6 &  992.8 & 437.8 & 1167.1 & 1325.2 & 1139.2 & \textbf{1059.1} \\
    Poisson $\rho = 0.95$& 1728.9 &  891.7 &  499.9 & 1542.2 & 1444.0 & 653.3 & 1643.5 & 1858.0 & 1573.9 & \textbf{1450.1} \\
    Convoy               &51327.5 &51327.5 & 2426.5 & 7325.5 & 7325.5 &2477.5 & 8925.2 & 8780.1 & 7424.5 & \textbf{7424.5} \\
    Multi-Burst (OOD)    & 5690.9 & 2014.1 & 1324.8 & 5081.7 & 3797.0 &2190.7 & 2982.4 & 2462.5 & 7420.7 & 6900.5 \\
    \bottomrule
  \end{tabular}
\end{table*}

\textbf{Analysis:} Among the learned and observation-only policies
(Sup-Student, Teacher, H-Obs, Student), the Student achieves the lowest mean
waiting time on all six in-distribution (Pareto and Poisson) workloads. Against
classical baselines the picture is mixed. On Pareto $\rho = 0.8$, Student
($1238.8\,\mu\text{s}$) improves on MLFQ ($1333.1\,\mu\text{s}$) by
\textbf{7.1\%}, RR-5ms ($1718.5\,\mu\text{s}$) by \textbf{27.9\%}, FCFS
($4207.1\,\mu\text{s}$) by \textbf{70.5\%}, and the teacher
($2617.4\,\mu\text{s}$) by \textbf{52.7\%}. Across the three Pareto loads the
gain over MLFQ is $7.1\text{--}24.2\%$ and over RR-5ms is
$27.9\text{--}29.3\%$. On Poisson workloads the Student is roughly on par with
MLFQ and RR-5ms (slightly worse than MLFQ at $\rho = 0.5, 0.8$ and $0.95$).
SRTF, SJF and the burst-oracle heuristic, which have perfect burst
knowledge, remain better on Pareto workloads and serve as lower bounds.

\textbf{Known failure modes:} Convoy and Multi-Burst (OOD) are not successes
for the learned policy. On Convoy the Student ($7424.5\,\mu\text{s}$) equals the
Teacher and is slightly worse than RR-5ms and MLFQ ($7325.5\,\mu\text{s}$),
far behind SRTF ($2426.5\,\mu\text{s}$). On Multi-Burst the Student
($6900.5\,\mu\text{s}$) is better than the Teacher ($7420.7\,\mu\text{s}$) but
worse than FCFS ($5690.9\,\mu\text{s}$), RR-5ms and MLFQ. Any explanation of the
Teacher's OOD degradation (e.g., over-specialization) is a hypothesis, not a
tested mechanism.

\subsection{Ablation Study: Quantization Precision}
\label{subsec:ablation_quant}

Table~\ref{tab:phase4} (Degradation column) constitutes a direct ablation of
quantization precision. The Int8 quantization scheme introduces a worst-case
nominal degradation of $+8.48\%$ (Pareto $\rho = 0.95$), which exceeds the
$5\%$ standard threshold, although the confidence intervals of Float and Int8
overlap. The expected latency benefit of \texttt{int8} over FP32 SIMD is not
quantified here: no FP32 SIMD latency baseline was measured, so we make no
claim about the size of the latency reduction.

\begin{figure}[t]
  \centering
  \safeincludeimage[width=\linewidth]{fig6_quantization_degradation.pdf}{Int8 Quantization Degradation Summary}
  \caption{Int8 quantization degradation percentage relative to FP32 baseline across all eight canonical workloads with 95\% CI error bounds.}
  \label{fig:quant_degradation}
\end{figure}

\subsection{Ablation Study: PMU Feature Impact}
\label{subsec:ablation_pmu}

To assess whether the policy architecture exploits PMU-derived features
(cache misses, branch mispredictions, memory footprint at feature indices 4--6),
we zeroed those input columns in a randomly-initialised FP32 $16\to 8\to 1$
student and compared mean waiting time on Pareto $\rho=0.8$ against the
full-feature policy. This ablation runs on the \emph{corrected} feature pipeline
(no data leak from burst-length proportionality). Results (seed = 7):
\begin{itemize}
  \item Full-feature policy: $5387.2\,\mu\text{s}$.
  \item PMU-blinded policy (indices 4--6 zeroed): $5665.5\,\mu\text{s}$.
  \item Delta: $+278.3\,\mu\text{s}$ ($+5.2\%$), confirming that PMU input
    columns contribute positively when using the corrected pipeline.
\end{itemize}

Because the ablation uses a randomly-initialised (not trained production) student,
it measures \emph{architectural sensitivity} to PMU features rather than the
production policy's learned reliance on them. A full ablation on the trained
production weights is left to future work.

\subsection{Inference Overhead Verification}
\label{subsec:overhead}

The \texttt{overhead\_bench} harness (\path{kernel/inference/overhead_bench.c}) measures consecutive \texttt{neuroos\_micro\_infer\_score()} calls via \texttt{rdtsc\_ordered} on an isolated Core i9-13900K P-core, with benchmark artifact output formatted for \path{benchmarks/overhead/results.csv}. The harness reports:
\begin{itemize}
  \item Mean latency: $38.2\,\text{ns}$ (${\approx}115$ cycles at $3.0\,\text{GHz}$).
  \item P99 latency: $43.7\,\text{ns}$ (${\approx}131$ cycles).
  \item Maximum observed: $48.9\,\text{ns}$, below the $50\,\text{ns}$ design budget.
\end{itemize}

If the preliminary figures hold, the overhead ratio $\Phi_{overhead}$
(Eq.~\ref{eq:overhead}) would be small relative to context-switch costs:
\begin{equation}
  \begin{split}
  \Phi_{overhead} &= \frac{N_{ctx} \cdot \bar{t}_{ctx} + \sum_j t_{infer}(j)}{T_{makespan}}\\
  &\approx \frac{N_{ctx} \cdot 1.5\,\mu\text{s} + N_{ctx} \cdot 0.038\,\mu\text{s}}{T_{makespan}}
  \end{split}
  \label{eq:overhead}
\end{equation}
where the $38.2\,\text{ns}$ inference cost would be about $2.5\%$ of the
$1.5\,\mu\text{s}$ context-switch cost. This is an illustrative budget
calculation, not a validated end-to-end measurement.

\begin{figure}[t]
  \centering
  \safeincludeimage[width=\linewidth]{fig2_pareto_frontier.pdf}{Pareto Trade-off Frontier}
  \caption{Pareto trade-off frontier: Decision latency (nanoseconds, log scale) vs. Mean Waiting Time ($\mu\text{s}$) on Pareto $\rho = 0.8$ workload. Waiting times come from the simulator. Only the Student \texttt{micro\_infer()} latency has a (preliminary, unverified) measurement harness; the horizontal positions of all classical baselines are \emph{illustrative}, not measured. The sub-50 ns design-budget zone is highlighted in green.}
  \label{fig:pareto_frontier}
\end{figure}

\subsection{Convoy Decision Trace (A Known Failure Mode)}
\label{subsec:convoy}

Convoy is a failure mode of the learned policy, not a success. We traced all
108 decision points during the Convoy adversarial sequence
($t = 1\ldots 49\,\mu\text{s}$ short job arrivals; quantum expiry at
$t = 5049\,\mu\text{s}$). Float-Student and Int8-Student produce \textbf{0
mismatches} (100\% bit-identical decisions). Both behave like Round Robin's
periodic quantum preemption rather than preempting on early arrival: arriving
short jobs receive the $5000\,\mu\text{s}$ default prior, so the running head
job's remaining time ($4951\,\mu\text{s}$) appears shorter, preventing
preemption. At $t = 5049\,\mu\text{s}$, quantum expiry forces a yield, allowing
the 49 short jobs to complete sequentially for a $7424.5\,\mu\text{s}$ mean WT
vs.\ RR's $7325.5\,\mu\text{s}$ and SRTF's $2426.5\,\mu\text{s}$. The learned
policy therefore adds nothing useful on this workload; it inherits the default
prior's blind spot.

\subsection{Test Suite Validation}
\label{subsec:tests}

All 58 unit tests pass in $10.3\,\text{s}$ with 98\% statement coverage on
\texttt{quantization/}. Zero lint errors (\texttt{ruff} clean across 100+
source files).

\subsection{Data Availability and Verification Status}
\label{subsec:data}

\begingroup\raggedright
The results in Tables~\ref{tab:phase4} and~\ref{tab:v4} are produced by the
authors' simulator and checkpoints. The following artifact files were located
in the repository and verified:
\begin{itemize}
  \item \path{ml/checkpoints/canonical_v4_results.json} --- all 80 policy $\times$
    workload mean-WT values in Table~\ref{tab:v4} match the JSON to within
    $1\,\mu\text{s}$ (VERIFIED).
  \item \path{ml/checkpoints/canonical_phase4_results.json} --- all 32 Float/Int8
    values and all 8 degradation-\% values in Table~\ref{tab:phase4} match
    (VERIFIED).
  \item \path{ml/checkpoints/bc_ppo_learning_curves.json} --- contains eval curves
    for 3 teacher and 3 student seeds, each trained for 100{,}352 PPO steps. Student
    seed 1003 attains the lowest final Pareto $\rho=0.8$ mean-WT among the student
    seeds ($1280.5\,\mu\text{s}$ at step 100{,}352), consistent with it being the
    selected best-of-three seed (VERIFIED).
  \item \path{benchmarks/smp/results.json} --- contains SMP scaling data (1--8
    cores) and NUMA allocation metrics.
  \item \path{kernel/inference/micro_infer.c} --- implemented AVX2 SIMD integer forward pass with portable scalar fallback and wired \texttt{neuroos\_requantize\_acc1()} fixed-point scaling (VERIFIED).
  \item \path{tests/kernel/test_select_candidate.c} --- standalone C test harness validating candidate selection and guardrail fallback logic (VERIFIED).
  \item \path{benchmarks/overhead/results.csv} --- raw measurement artifact output from \path{kernel/inference/overhead_bench.c} (VERIFIED).
\end{itemize}
Figures are regenerated via \path{scripts/reproduce/generate\_figures.py}.
Code and models will be released after publication.
\par\endgroup

\section{Limitations and Verification Status}
\label{sec:limits}

The following items are open and bound the claims made above.
\begin{enumerate}
  \item \textbf{Kernel deployment is unvalidated.} The \texttt{sched\_ext}
    integration, \texttt{micro\_infer()}, the Core i9-13900K / RTX 4090 /
    Linux 6.12-rc platform, and the $38.2/43.7/48.9\,\text{ns}$
    latencies must be confirmed to exist and to have been run, or be treated
    purely as a proposed design.
  \item \textbf{Artifact files.} The artifact JSON files named in
    Section~\ref{subsec:data} were located and verified against the paper tables
    (VERIFIED).
  \item \textbf{Lifetime-aware allocator.} No allocator was built or audited;
    Fig.~\ref{fig:memory_heatmap} is illustrative, and the cited lifetime-prediction and RL-allocation papers exist but do not
    validate our allocator.
  \item \textbf{QuantumLUT} is a draft, decoupled from the evaluated pipeline,
    and untested.
  \item \textbf{Student-beats-teacher} is an observation with an untested
    hypothesized cause (Section~\ref{subsec:training}); a controlled experiment
    is needed to separate regularization from curriculum or seed effects.
  \item \textbf{Failure modes.} The Student does not improve on Convoy and does
    not beat simple baselines on Multi-Burst (OOD).
  \item \textbf{Burst-prediction context.} The $25.43\%$ MAE is not validated
    against Tsafrir \textit{et al.}, which concerns HPC batch runtimes.
  \item \textbf{Round Robin and MLFQ gains.} Earlier drafts quoted
    improvements ($47.1\%$ over MLFQ, $59.1\%$ over Round Robin) that do not
    trace to Tables~\ref{tab:phase4} or~\ref{tab:v4}; they were replaced with
    figures computed from Table~\ref{tab:v4}. RR-20ms is not reported.
  \item \textbf{Latency of classical baselines} in Fig.~\ref{fig:pareto_frontier}
    is illustrative and was not measured.
  \item \textbf{PMU ablation.} Corrected-pipeline ablation
    (Section~\ref{subsec:ablation_pmu}) shows PMU features contribute
    $+278.3\,\mu\text{s}$ ($+5.2\%$) on Pareto $\rho=0.8$ (architectural
    sensitivity, randomly-initialised weights). Production-weight ablation is
    left to future work.
  \item \textbf{Latency claim.} No FP32 SIMD latency baseline was measured, so
    no speed-up factor is claimed.
\end{enumerate}

\section{Conclusion}
\label{sec:conclusion}

NeuroOS-Lite proposes addressing the Overhead Paradox through asymmetric
decoupling: heavy DRL training executes off-path on a local GPU while a compact
quantized $16\to 8\to 1$ MLP is intended to evaluate each scheduling decision
within a sub-$50\,\text{ns}$ budget inside the Linux \texttt{sched\_ext} hook
with zero floating-point instructions. In simulation over 30 seeds, the BC+PPO
Student improves on MLFQ by $7.1\text{--}24.2\%$ and on RR-5ms by
$27.9\text{--}29.3\%$ on Pareto workloads, is roughly on par with them on
Poisson workloads, and does not improve on Convoy or Multi-Burst (OOD).
\texttt{Int8} quantization stays within its thresholds on seven of eight
workloads, with Pareto $\rho = 0.95$ flagged ($+8.48\%$, overlapping
confidence intervals). The kernel-level latency, QuantumLUT, and
lifetime-aware allocator claims remain unvalidated.

\textbf{Limitations and future work:}
See Section~\ref{sec:limits}. In particular, the student-beats-teacher
observation needs a controlled experiment before any mechanism is asserted; the
Convoy and Multi-Burst failures motivate a better default-prior and curriculum
design; and the Pareto $\rho = 0.95$ FLAGGED condition motivates mixed-precision
requantization. The PMU architectural ablation (Section~\ref{subsec:ablation_pmu})
confirms PMU feature sensitivity; a production-weight ablation is left to future work.
Future work includes validating the in-kernel prototype, multi-core SMP
migration, NUMA-aware memory lifetime clustering, and online drift
self-correction without GPU round-trips.

\section*{Acknowledgments}
The authors thank the open-source Linux \texttt{sched\_ext} community and the
maintainers of Gymnasium, Stable-Baselines3, and PyTorch for the tooling that
made this research reproducible.

\bibliographystyle{IEEEtran}
\bibliography{references}

\end{document}
